\section{Further research questions and proposed approaches}
\label{sec:research-program}

The research priorities below are chosen to connect the proved local
results to a complete hierarchy.  Each has a concrete success criterion.
Several are mathematical questions rather than requests for more
benchmarking; a fast implementation cannot substitute for their missing
proofs.

\subsection{Priority I: preserve geometry while shrinking the state}

\begin{problem}[Patterned ball absorption with shared source addresses]
\label{prob:patterned-balls}
Given a verified unused-cocore decomposition of an irreducible manifold,
a boundary pattern in explicit or compressed source charts, and the
rooted forest of Theorem~\ref{thm:attachment-kernel}, construct a
pattern-preserving reduced presentation and a verifiable transport map.
Bound its encoded size in terms of the source size, the retained genus
budget, and a precisely defined pattern-interface parameter.
\end{problem}

The promising approach is to retain an absorbed ball only when its pattern
has information that cannot be carried through its parent collar by the
chosen chart representation.  The other balls become shared transport
nodes.  Proposition~\ref{prop:forest-transport} prevents the forest paths
themselves from causing quadratic duplication, once the local transport
maps have been constructed.  A successful proof must also control the
complexity of those maps.  Merely storing one Boolean ``marked'' flag per
ball does not encode an embedded graph on its boundary, and a fixed genus
does not bound the normal weight of a curve on that boundary.

A useful first target is a pattern that is a disjoint family of arcs,
with every endpoint on an explicitly listed attaching disc.  The next
target allows trivalent vertices and closed components.  For each class,
one should either prove a polynomial compressed transport bound or exhibit
a family forcing the proposed representation to grow.  The theorem should
state which isotopies are allowed at the retained core and which boundary
components remain individually labelled.  These are part of the input
contract, not implementation details that can be chosen after a proof.

\begin{problem}[A normal-coordinate adapter for the retained core]
\label{prob:core-normal-adapter}
Turn the initial admissible handle subpresentation into a normal or
regulated-surface coordinate system whose solutions have an explicit
source interpretation.  Construct a polynomial-time map for the surfaces
and boundary data actually needed by the next hierarchy step.
\end{problem}

The four local characteristic graphs give a finite starting catalogue.
For each graph, enumerate the allowed local disc types, their boundary
arcs, and matching equations.  Prove that assembling these records gives
exactly the intended embedded surfaces, including the restrictions needed
to avoid incompatible local choices.  The main challenge is the relation
between this inherited handle structure and any new triangulation: a
small count of local types is not, by itself, a bound on every global
coordinate or a surface-search algorithm.  An adapter tested on the
canonical trefoil and figure-eight cores would be a concrete first
implementation milestone, but its completeness needs a separate proof.

\begin{problem}[Reconstruction relative to a specified pattern]
Strengthen Corollary~\ref{cor:normal-cut-reconstruction} to an output
that retains a specified boundary pattern, seam identification and
normal-coordinate charts.  State and prove the additional parameter
needed beyond $t+b_0(S)+\chi_-(S)$.
\end{problem}

There are two useful output models to compare.  One emits an ordinary
triangulation with the marking expanded into its individual normal arcs;
the other retains a small triangulation and a compressed marking or
transport circuit.  A binary normal-coordinate vector is itself a
compressed representation and must be distinguished from the first model.
The research should begin by finding families that separate precisely
specified encodings, for example repeated twists in a boundary collar.
Large normal weight alone does not imply large binary coordinate length.
The extension of a vertical-annulus twist used in
the unmarked proof explains why the unmarked complexity can stay small
while a prescribed marking changes.  Algorithms for curves in surfaces
\cite{Lackenby2026Curves} are candidates for comparing or normalizing the
resulting compressed markings, once the three-dimensional chart transfer
has been supplied.

\subsection{Priority II: make native orbit compression quantitative}

\begin{problem}[A uniform trace bound for ownership-weight queries]
\label{prob:weight-runs}
For the particular maintained weighted-orbit implementation used here,
prove a polynomial upper bound on its intermediate weight runs, cycles
and replay-certificate size for authenticated native ownership queries,
as a function of source records and binary interval lengths.
\end{problem}

The published weighted orbit theorem gives a polynomial algorithm for
the abstract task.  It does not automatically give such a bound for a
different implementation's exact scheduling and proof trace.  The new
length field adds owner intervals that were absent from the sparse
endpoint-mask input.  This raises the observed run count on some cases
even while removing a whole census.  A correct proof must therefore
track owner-interval fragmentation, not reuse a port-only sparse-token
bound.

One possible route is to decompose the weight function into two channels:
a small terminal channel and a structured edge-mass channel.  The
terminal channel has a conservation certificate; the mass channel is
the image of disjoint affine owners.  If native geometric constraints
force a bounded number of mass breakpoints per surviving orbit phase,
they may give a sharper bound than an arbitrary weighted-query analysis.
If this invariant fails, a counterexample should guide a revised
representation or a direct implementation of the published weighted
algorithm.  Maintaining exact independent replay is a useful requirement
for either route.

\begin{problem}[Optimal deterministic pair labels]
Let $r$ labels be inserted once each into an unknown partition.  Find
nonnegative integer labels with the smallest maximum bit length that
permit exact unordered-pair decoding, a separate additive length channel,
and a verifiable global inventory, without a quadratic decoding table.
\end{problem}

The pair code in Section~\ref{sec:native-moments} uses a quadratic moment
and a radix of order $r^3$.  It has $O(\log r)$-bit labels and an explicit
integer-square-root decoder.  For a codec that must recover each pair
from its scalar sum alone, the $\binom r2$ distinct unordered pairs must
have distinct sums.  Labels drawn from $[0,M]$ then give only $2M+1$
possible sums, so
\[
  2M+1\geq\binom r2,\qquad M=\Omega(r^2).
\]
This lower bound applies to that local additive decoding contract; it is
not a lower bound for every encoding that might exploit the complete
inventory globally.  It leaves a factor of order $r$ between the present
magnitude and the counting bound, while the bit-length order is already
optimal for the local contract.  Closing this gap requires both smaller
values and an efficient inverse for their pair sums.

There is a useful way to separate the inventory requirement from that
inverse problem.  Suppose $u_i\in[0,M]$ have distinct sums on distinct
unordered pairs.  Shift them to $v_i=M+1+u_i$.  Every singleton is then at
most $2M+1$, whereas every pair sum is at least $2M+2$.  Pair-sum
injectivity is unchanged.  With an owner radix $D>2\max_i v_i$, zero-
and one-label orbits cannot masquerade as a decodable pair, and the
consumption argument of Lemma~\ref{lem:native-pair-global-inventory}
applies under the same complete-histogram and multiplicity checks when
the recovered matching exhausts the authenticated inventory.  Thus an
efficiently invertible Sidon-type construction would
be a candidate; existence of distinct pair sums alone is insufficient.
Any change should be benchmarked against the current power-of-two radix,
whose divisions are bit shifts.

\begin{problem}[Compatible boundary frames and passive transport]
Augment the native boundary output with local fibre frames on all relevant
base-boundary components, including passive components, together with the
transition and orientation-obstruction data required to compare them.
Prove that these data suffice for a stated gluing or surgery operation.
No global fibre orientation is required when the bundle obstructs one.
\end{problem}

The present output gives local fibre transport across active gaps and
retains passive multiplicities through the optional full packet.  It
does not supply a chosen parametrisation of every passive boundary circle
or source paths comparing all local frames.  For an orientable total
interval bundle, fibre holonomy around a base-boundary circle is already
trivial: its vertical boundary is an annulus, not a M\"obius band.  Thus
the missing data here are parametrisations and source-based frame
comparisons, rather than an undetermined boundary fibre flip.  A
finite-group lift from the incoming moment report is a candidate primitive
for transport along specified source paths, including interior paths used
to compare frames.  It must be checked against the bundle's orientability
data.
Testing a few nonorientable bases is useful for finding mistakes; the
main result should be a sufficiency theorem for the complete gluing
record.  The certificate must preserve components with no active port,
because a terminal-only census cannot see them.

\subsection{Priority III: bound a complete hierarchy}

\begin{problem}[A complete quantitative hierarchy theorem]
\label{prob:global-depth}
Prove, or locate and verify in full detail, the precise hierarchy
construction needed for a bound of $L=O((\log n)^2)$ together with the
surface and pattern parameter $G=n^{O(1)}$ and polynomial local work.
Then identify the additional argument required for the stronger
$L=O(\log n)$ version of the announcement.
\end{problem}

The February and March 2021 statements should be tracked separately.
The 2026 preprint is a substantial source of rigorous machinery, but
its iteration estimate does not furnish every bound in
Theorem~\ref{thm:conditional-accounting}.  A useful research artifact
would map each assumption of that theorem to a proved proposition,
including its actual parameter and whether it applies with empty,
essential, or arbitrary boundary pattern.  This would expose any missing
hypothesis before implementing a search that relies on it.

\begin{problem}[Discover useful surfaces without exhaustive enumeration]
Develop a complete method to select the normal or regulated surfaces
needed by that hierarchy within the target total budget.  Bound the
encoded normal coordinates and the number of unsuccessful candidates.
\end{problem}

The new initial core removes certified irrelevant source handles, so
it is a better input to such a search.  The genus kernel may also allow
the search to focus on a small number of interacting nonball pieces.
Nevertheless, a single retained torus-boundary piece can already be
an arbitrary knot exterior.  Neither reduction permits a claim that
enumerating all fundamental surfaces becomes quasipolynomial.
A promising direction is a selector tied to a proved hierarchy potential,
with a certificate that either supplies the required surface or justifies
one of the special-region alternatives.  A heuristic ranking can help
experiments, but a worst-case proof must bound the fallback search too.

\begin{problem}[Control intersections in multisurface states]
Give a compressed, source-authenticated representation of the intersecting
surfaces permitted by the announced multisurface method, and prove that
its size remains within the state bound of
Theorem~\ref{thm:conditional-accounting}.
\end{problem}

Proposition~\ref{prop:isotropic-obstruction} rules out an easy replacement
by a disjoint basis.  The relevant interface must encode intersection
curves and arcs, their order on source pieces, and the effect of cutting
in successive complements.  Native first-return queries may recover
some of the necessary transitions.  Sparse cut grammars and disconnected
tile kernels from the incoming packages may then compress repeated
regions.  The research burden is to prove both geometric realizability
and a bound on the width or number of such regions along the actual
hierarchy.  A fast grammar algorithm on unrestricted input does not
provide that geometric bound.

\begin{problem}[A combined potential for topology, interfaces and marks]
Find a potential that controls hierarchy depth while charging retained
nonball deficits, attachment cycles, active interfaces, and the cost of
the boundary pattern in one state representation.
\end{problem}

The exact identity $\delta=\beta+\sum(1-\chi(K_v))$ supplies an additive
topological budget.  The moment compiler supplies an output bound in the
number of active ports and logarithmic lengths.  These fit naturally
into a proposed potential with separate topological and interface terms.
The difficult issue is monotonicity: a valid cut can increase boundary
genus and expose new ports.  A successful potential must prove that such
increases are paid for by progress in the global hierarchy, or by a
special-region reduction.  Simply adding the available size parameters
does not prove any decrease.  This is a central mathematical question,
not a matter of choosing convenient benchmark weights.

\subsection{Priority IV: make performance evidence useful for integration}

\begin{problem}[Choose source operators by measured structural cost]
Determine when the maintained boundary-period certificate, the unused
component certificate, the four-field moment codec, or the compact pair
codec has the best completed production-and-replay cost.  Design a
deterministic, inexpensive selection rule with a correctness-preserving
fallback.
\end{problem}

The delivered results already rule out an unconditional wall-clock claim.
The unused-handle certificate has extra component records and can be
larger on small sources; the ownership query can generate more weight
runs; padded trefoil examples may become slower despite a smaller core.
Candidate predictors include surviving attaching occurrences, the number
of zero-degree handles, owner-run count, port count, and binary coordinate
length.  They should be obtained from work already required by the source
compiler whenever possible.  The fallback must retain exact status and
certificate semantics; a timeout is not evidence of knottedness.

\begin{problem}[Extend independent evidence beyond the present fixtures]
Build a stable corpus of valid relabelled triangulations, reducible
controls, high-valence source edges, large binary normal coordinates,
port-free bundles, and independently known knot types.  Measure both
local operations and end-to-end recognition when an opt-in integration
exists.
\end{problem}

The immediate deliverable should be a source-identified record for every
case, its expected theorem domain, exact proof replay, and raw samples.
Separate favorable asymptotic families from finite stress tests.  Retain
negative cases and historical baselines.  The present package follows
that convention, but its finite banks are not a random sample of all
knot exteriors or all normal surfaces.  A broad integration evaluation
should follow the mathematical and representation milestones above,
so that any observed speedup can be attributed to a completed operation.

\subsection{A proposed order of work}

The most direct next milestone is the normal-coordinate adapter on the
four initial local templates, together with an opt-in call to the
authenticated initial compiler.  In parallel, the ownership-weight
trace bound can be pursued without changing the topology.  The next
geometric milestone is patterned ball absorption and compatible bundle
charts.  Only then is it reasonable to combine the incoming grammar and
tile kernels into a recurring hierarchy state.  The global depth and
surface-selection arguments should be developed alongside that work,
with every needed quantitative claim recorded explicitly.  This ordering
makes each implementation step consume a proved interface and makes
failure to achieve the target exponent mathematically diagnosable.
